\chapter{Ada for C Programmers}
\label{ch:ada-primer}

The rest of this book assumes you can read Ada. If you are coming from C or C++ --
as most embedded engineers are -- this chapter is the on-ramp: enough of the
language to follow every example that follows, framed against the C you already
know. It is not a full Ada course (for that, see the references in
Appendix~\ref{ch:references}); it is the slice this book leans on, in one sitting.

The recurring theme is that Ada says the same things C does but makes more of them
\emph{checkable}: types are distinct, the direction of every parameter is declared,
and visibility is enforced by the compiler rather than trusted to a header.

\section{Declarations and the three ``equals''}

A declaration names something and gives its type; an initial value is optional for
a variable and required for a constant:

\begin{lstlisting}
   Count : Integer := 0;            -- a variable, initialised
   Limit : constant Integer := 8;   -- a constant (it can never be assigned to)
   Pin   : ESP32S3.GPIO.Pin_Id;     -- a variable, no value yet
\end{lstlisting}

Three symbols a C programmer reads as ``equals'' are kept distinct in Ada, and the
distinction is worth fixing now because all three turn up on the same page later:

\begin{itemize}
  \item \texttt{:=} is \textbf{assignment} -- ``becomes''. \texttt{Count := 5;}
  \item \texttt{=} (and its negation \texttt{/=}) is \textbf{equality} -- a Boolean
        test, as in \texttt{if Count = 5}. It never assigns.
  \item \texttt{=>} is \textbf{association} -- ``goes with''. It names which
        parameter, array element, or alternative a value belongs to, as in
        \texttt{Configure (Pin, Mode => Output)}.
\end{itemize}

There is no \texttt{=} that assigns and no \texttt{==}; the classic C slip of
writing \texttt{if (x = 5)} when you meant \texttt{==} cannot even be expressed.

\section{Subprograms and parameter modes}
\index{parameter modes}\index{subprogram}

A \emph{procedure} performs an action; a \emph{function} returns a value. Both take
parameters, and -- unlike C -- each parameter declares its \emph{direction}:

\begin{lstlisting}
   procedure Read (Reg : in  Register_Id;   -- caller -> subprogram (the default)
                   Val : out Word);         -- subprogram -> caller

   procedure Update (Count : in out Natural);  -- both ways

   function Is_Set (W : Word; Bit : Natural) return Boolean;
\end{lstlisting}

\begin{itemize}
  \item \texttt{in} (the default) passes a value the subprogram may read but not
        change -- a C value parameter.
  \item \texttt{out} is how a procedure \emph{returns} data: the caller passes a
        variable and the subprogram fills it. This is why so many operations in this
        book end in \texttt{Result : out Status} or \texttt{Val : out Word} -- it is
        the idiomatic way to hand back several results at once without inventing a
        struct.
  \item \texttt{in out} both reads and updates the caller's variable.
\end{itemize}

In C you imitate \texttt{out} and \texttt{in out} with pointers (\texttt{int *val})
and trust the caller to know which the function reads and which it writes. Ada makes
the direction part of the declaration, so the compiler -- and the reader -- always
know. Arguments may be given by position or by name, and naming is preferred where
a bare value would be a riddle:

\begin{lstlisting}
   Configure (Pin => 2, Mode => Output);    -- named: self-documenting
   Configure (2, Output);                   -- positional: the same call
\end{lstlisting}

A parameter may have a default (\texttt{Port : Port\_Type := 502}), so a caller can
omit it -- which is why \texttt{Connect (S, "10.0.0.7")} compiles even though
\texttt{Connect} has half a dozen parameters.

\cnote{a function with several results returns a struct or scribbles through
pointers, and nothing in the signature says which pointer is in and which is out.
The \texttt{in}/\texttt{out}/\texttt{in out} mode is that missing information, moved
into the declaration and checked.}

\section{Packages: the unit of modularity}
\index{package}

A \emph{package} is Ada's module, and it comes in two parts, usually two files:

\begin{itemize}
  \item the \textbf{specification} (\texttt{.ads}) -- the visible interface, like a
        C header: the types and subprograms a client may use;
  \item the \textbf{body} (\texttt{.adb}) -- the implementation, like the \texttt{.c}.
\end{itemize}

The difference from C is that the split is \emph{enforced}, not conventional. A
client sees only what the specification declares; anything in the body, or below a
\texttt{private} line in the spec, is genuinely inaccessible -- the compiler refuses
it, rather than the linker merely happening not to find it.

Two clauses bring a package into view:

\begin{lstlisting}
   with ESP32S3.GPIO;     -- make the package available (like #include, but typed)
   use  ESP32S3.GPIO;     -- AND open its names, so the prefix can be dropped

   ESP32S3.GPIO.Toggle (Led);  -- with only
   Toggle (Led);               -- with + use
\end{lstlisting}

\texttt{with} is the dependency: it says ``this unit needs that one,'' and the
compiler tracks it -- there is no include order to get wrong and no header guards.
\texttt{use} is optional convenience that lets you drop the package prefix. This
book \texttt{with}s widely and \texttt{use}s sparingly, so most calls keep their
prefix and you can see at a glance where each name comes from.

Packages nest. \texttt{ESP32S3.W5500.Sockets} is a \emph{child} of
\texttt{ESP32S3.W5500}, itself a child of \texttt{ESP32S3}: the dotted name is a
hierarchy, and that is how the HAL is laid out -- a chip namespace with one package
per peripheral. This is enough to read the code; the next chapter
(Chapter~\ref{ch:packages}) makes the case that the package is the load-bearing idea
of the whole design -- reuse, privacy, naming and the build all in one construct.

\section{Control flow}

The control structures are the C ones with words in place of braces, plus a couple
of conveniences:

\begin{lstlisting}
   if Count = 0 then
      ...
   elsif Count < Limit then
      ...
   else
      ...
   end if;

   case Status is
      when OK        => ...    -- no fall-through, so no 'break'
      when Timed_Out => ...
      when others    => ...    -- everything not named above
   end case;

   for I in 0 .. 7 loop ... end loop;   -- I is declared by the loop, read-only
   while not Done loop ... end loop;
   loop                                  -- a bare loop ...
      exit when Last < First;            -- ... left by an explicit condition
   end loop;
\end{lstlisting}

The differences that catch a C programmer: \texttt{case} alternatives never fall
through (there is no \texttt{break}) and must cover every value (hence
\texttt{when others}); and a \texttt{for} variable is declared by the loop, ranges
over a type or range, and cannot be assigned inside the body.

\section{Types, attributes, and aggregates}

Ada's typing is \emph{strong}: two distinct types never mix without an explicit
conversion, even when both are integers -- a channel count and a clock divider
should not be interchangeable just because both are numbers. The everyday integer
types are \texttt{Integer}, \texttt{Natural} ($\geq 0$) and \texttt{Positive}
($\geq 1$); the embedded-specific ones -- range-constrained subtypes and modular
(wrap-around) types for registers -- are the subject of Chapter~\ref{ch:embedded}.

\textbf{Attributes} are built-in queries written with an apostrophe, the ``tick.''
They have no C equivalent and appear on nearly every page:

\begin{lstlisting}
   A'First   A'Last   A'Range   A'Length    -- about an array or a range
   T'Image (X)                               -- a value rendered as text
   Obj'Access                                -- a typed, checked pointer to Obj
\end{lstlisting}

\textbf{Arrays} carry their own bounds, and the bounds may start anywhere:

\begin{lstlisting}
   type Byte_Array is array (Natural range <>) of Byte;  -- size set per object
   Buf : Byte_Array (0 .. 63);                            -- this one is 0 .. 63
\end{lstlisting}

\textbf{Aggregates} build a whole composite value at once -- by position, by name,
or with \texttt{others} for the rest:

\begin{lstlisting}
   P : Point      := (X => 1, Y => 2);          -- a record, by name
   Z : Byte_Array (0 .. 3) := (others => 0);    -- every element set to 0
   M : MAC_Address := (16#48#, 16#CA#, 16#43#, 16#37#, 16#CB#, 16#5C#);
\end{lstlisting}

That \texttt{=>} is the same one as in a named call: it ties a value to the
component it fills.

\section{Objects: tagged types and dispatching}
\index{tagged type}\index{dispatching}\index{interface}

Ada's object orientation is built on the \emph{tagged} record -- a record you are
allowed to extend. A C++ programmer can read the next few lines as classes and
virtual methods; the words differ, the machinery does not.

\begin{lstlisting}
   type Server is tagged record               -- a base "class"
      ...
   end record;
   procedure Handle (Self : in out Server; ...);   -- a "primitive operation"

   type My_Server is new Server with record    -- a derived "class"
      Data : Word_Array (0 .. 63);             -- adds its own fields
   end record;
   overriding procedure Handle (Self : in out My_Server; ...);  -- overrides it
\end{lstlisting}

A subprogram that takes the type as its first parameter is a \emph{primitive
operation} -- a method -- and a derived type inherits the primitives it does not
\texttt{override}. The keyword \texttt{overriding} is optional but wise: written, it
makes the compiler confirm you really are replacing an inherited operation and not
quietly adding a new one through a typo.

\emph{Dispatching} is virtual-call behaviour, and it turns on the class-wide type
\texttt{T'Class} -- ``\texttt{Server} or anything derived from it.'' A call through a
\texttt{Server'Class} value runs the \emph{actual} object's version:

\begin{lstlisting}
   procedure Drive (S : in out Server'Class) is
   begin
      Handle (S, ...);   -- dispatches: runs My_Server.Handle if S is a My_Server
   end Drive;
\end{lstlisting}

An \emph{interface} is the limiting case: a tagged type with operations but no data,
all of them to be supplied by whoever implements it -- C++'s pure-abstract class.
This is exactly the shape of the book's pluggable backends. \texttt{Net\_Devices.%
Device} is an interface every network chip implements; \texttt{Modbus.Slave.Server}
is a tagged base you derive from, overriding only the request handlers your device
supports while the framework calls them by dispatch. You meet that pattern in
Part~IV; this is the machinery under it.

\section{A C-to-Ada cheat sheet}

\begin{table}[ht]
\centering
\begin{tabular}{@{}ll@{}}
\toprule
\textbf{C / C++} & \textbf{Ada} \\
\midrule
\texttt{\#define N 8}          & \texttt{N : constant := 8;} \\
\texttt{const int}             & \texttt{constant} \\
\texttt{int x = 0;}            & \texttt{X : Integer := 0;} \\
\texttt{x = 5;}                & \texttt{X := 5;} \\
\texttt{x == 5}                & \texttt{X = 5} \\
\texttt{x != 5}                & \texttt{X /= 5} \\
\texttt{struct} / \texttt{enum} & \texttt{record} / enumeration type \\
\texttt{typedef}              & \texttt{type} / \texttt{subtype} \\
\texttt{T *p}                 & \texttt{P : access T} \\
\texttt{\&x}                   & \texttt{X'Access} \\
\texttt{void f(int *out)}     & \texttt{procedure F (V : out T)} \\
\texttt{\#include "h.h"}       & \texttt{with H;} \\
header / \texttt{.c}          & spec (\texttt{.ads}) / body (\texttt{.adb}) \\
\texttt{for (;;) \{\}}         & \texttt{loop ... end loop;} \\
\texttt{switch}/\texttt{break} & \texttt{case ... when =>} (no break) \\
\texttt{0xFF}, \texttt{0b1010} & \texttt{16\#FF\#}, \texttt{2\#1010\#} \\
\texttt{x \& mask}            & \texttt{X and Mask} (on a modular type) \\
\bottomrule
\end{tabular}
\caption{The C constructs this book's Ada stands in for.}
\end{table}

That is the whole of the language you need to start reading. The features that make
Ada worth the trade -- subtypes that fence off illegal values, modular types for
registers, contracts, controlled types, and the tasking model -- are introduced
where the book first relies on them, beginning with Chapter~\ref{ch:embedded}.
